\section{Experiments}
\label{sec:protocol}
\subsection{Datasets}
We evaluate on three settings derived from LUNA16~\cite{setio2017,armato2011} and HiPaS~\cite{chu2025hipas}, whose fitting, development, calibration and test roles are listed in Table~S1 of the supplement. In each setting a \emph{block} is a patient (LUNA16) or a scan (HiPaS), and every control crop is paired with a bank of 18 insertions that vary lesion shape, contrast and offset.

The \emph{low-texture LUNA16 subset} comprises 549 neighborhoods from 183 patients, each with a single, crossing or skew added vessel. Its core intensity standard deviation is capped at $42\hu$ to isolate cases in which vessels compete strongly with the lesion, and its lesions are compact, radially symmetric signals. Thresholds and the analysis for this subset were fixed before testing, but because these cases also informed the design of the shared readouts, we treat readout comparisons on them as development results.

The \emph{broader LUNA16 held-out set} is the main test set. It uses 170 patients disjoint from all fitting, development and calibration data, with four parenchymal neighborhoods per patient (680 controls and 12,240 insertions). It imposes no texture cap (median core SD $117\hu$), uses added vessels that bend and taper, and introduces bilobed and part-solid-like lesions at three contrasts (84, 196 and $364\hu$) and three offsets (0, 4 and 7\,mm) from the proposal center. These shapes differ from the search templates, so the test does not reward template matching; the sampling constraints and the lesion model are detailed in Sec.~S3.

Finally, the \emph{HiPaS native-vessel set} tests whether the effect survives real anatomy. It consists of 72 noncontrast scans from a different source with manual artery and vein annotations. Six neighborhoods per scan are centered between nearby annotated arteries and veins, yielding 432 controls and 7,776 insertions without any added vessel (median core SD $253\hu$). The annotations define only the sampling population and the analysis strata; the scorer observes the image alone.

\subsection{Experimental Settings}
Classifiers for the broader set are fitted on 80 patients and tuned on 85 others over the three coupled settings of $(\rho,\tau)$ and the ridge strength (Sec.~S2). Before the test images were rendered, we fixed the cohort, the insertion bank, the analysis plan and three scorers selected on development: the aligned profile as the \emph{lead} scorer, fixed subtraction as the strongest alternative readout, and the rotated profile with a refitted classifier. The primary criterion was a crop-sensitivity gain of at least 5\,pp at 10\% FPR, with a positive lower confidence limit against both the rotation and fixed subtraction; localized sensitivity was a prespecified secondary endpoint. A development bootstrap indicated that approximately 123 patients provide 80\% power to detect a true 5-pp gain over the rotation, and 74 patients over fixed subtraction, so the 170 test patients are sufficient for both comparisons.

\subsection{Baselines}
\label{sec:baselines}
All readouts of Sec.~\ref{sec:readout} are evaluated with both placements, together with three reference scorers. The \emph{CT-only} scorer uses the shared features and the CT-only readout, and the \emph{CT + shell} scorer adds the four shell summaries, which supply the exterior information without a spatial field. The \emph{Hessian} baseline subtracts the same forecast at a strength of $0.35$, $0.7$ or $1.0$ before computing a standard compact-curvature statistic at Gaussian scales of 1.3, 2.5 and $4.0\mm$; it measures the contribution of a different readout when the nuisance estimate is held fixed. The learned baseline is a hybrid CNN with a small 3D feature extractor and a shared CT feature path, whose second input is either smoothed CT (the \emph{CT-only CNN}), the aligned forecast, the rotated forecast or the residual $y-f$; architecture, data, seeds and training budget are identical across inputs (Sec.~S6). Apart from the detector-score analysis on LNDb in Sec.~\ref{sec:calibration}, we do not compare against complete nodule-detection systems: they couple proposal generation, context and classification, so a placement sham cannot be applied to their vessel evidence alone, which is the quantity under test.

\subsection{Evaluation Metrics}
We report two operating points. The \emph{matched} operating point sets each scorer to a 10\% candidate false-positive rate (FPR) on the evaluation controls via the order statistic of their scores, and therefore measures ranking at a fixed budget. The \emph{frozen} operating point applies cutoffs derived from 700 independent negative crops of 175 LUNA16 calibration patients, accepting scores strictly above the 10\% order statistic; its achieved FPR may differ between scorers. At either operating point, \emph{crop sensitivity} is the fraction of lesion crops accepted, and \emph{localized sensitivity} additionally requires the returned location to lie within 3\,mm of the insertion center $x_\star$,
\begin{equation}
\Pr\bigl(A(y)=1,\ \|L(y)-x_\star\|\le 3\mm \mid Y=1\bigr).
\label{eq:locsens}
\end{equation}
Controls count against the FPR budget irrespective of their returned location.

Differences between two scorers are reported as ``$+x\pp\,[a,b]$'': $x$ is the paired difference in sensitivity in absolute percentage points (pp), and $[a,b]$ its 95\% confidence interval from 3,000 bootstrap resamples of blocks, stratified by release fold. Each resample retains a block's full insertion bank and controls, so both scorers are always evaluated on the same crops. For matched intervals the cutoff is recomputed in every draw. For frozen cutoffs we report intervals conditional on the calibration set and, separately, intervals that also resample the calibration patients.
